% Cover letter for the CMAME submission (13 Oct 2026).  Author to review, date and sign.
\documentclass[11pt]{letter}
\usepackage[margin=2.5cm]{geometry}
\usepackage[T1]{fontenc}
\signature{Kumar Kaushik Ranjan\\Department of Mechanical Engineering\\Tolani Maritime Institute, Pune, India\\kumarr@tmi.tolani.edu}
\address{Department of Mechanical Engineering\\Tolani Maritime Institute\\Induri, Talegaon Dabhade, Pune 410507\\Maharashtra, India}
\begin{document}
\begin{letter}{The Editors\\Computer Methods in Applied Mechanics and Engineering}
\opening{Dear Editors,}

I submit the manuscript \emph{``Operator-adapted learned multigrid cycles for high-contrast elliptic problems: accuracy,
failure modes and cost''} for consideration as a research article in \emph{Computer Methods in Applied Mechanics and
Engineering}.

The paper is a controlled methodological study. It asks what learned smoothing adds to one pass of a multigrid cycle for
high-contrast diffusion problems on unstructured finite element meshes, when the transfers, Galerkin coarse operators and
near-kernel deflation vectors are computed from each sample's operator and only the smoothers are learned. The architecture
provably represents the classical cycles with fixed polynomial smoothers on the same hierarchy, is positively homogeneous in
the load, and has an energy-norm safeguard. This allows a paired comparison with a classical cycle that receives exactly the
same hierarchy, coarse solution, deflation and output projection.

At four update stages per smoother, the learned pass is about ten times more accurate in the energy norm than the
degree-four Chebyshev-smoothed pass on conductive channels and about twice as accurate on insulating barriers, reproducibly
across held-out data seeds and training runs; ablations attribute parts of this gain to the operator-adapted transfers,
deflation, Galerkin output and nonlinearity. The paper reports with equal weight where the result stops: the margin shrinks
quickly when the classical pass receives more smoothing, and three classical passes are more accurate at a third of the
cost; a network trained with one load loses its advantage on other loads; floating islands keep substantial error; and the
implementation is 12--98 times slower than converged classical solvers at matched accuracy. The manuscript therefore does
not claim a fast solver or a broadly superior neural operator. Its contribution is a reproducible measurement of what
learned smoothing contributes inside a prescribed multilevel pass, and of why that contribution does not yet translate into
solver efficiency, which I believe is useful to readers designing learned components for multilevel solvers in
computational mechanics.

All tables are regenerated from stored per-sample records by scripts that, with the code, records and checkpoints, are
available from me on request. The manuscript has not been published and is not under consideration elsewhere. The use of generative
AI tools in preparing this work is declared in the manuscript. There are no competing interests and no specific funding.

\closing{Yours sincerely,}
\end{letter}
\end{document}
